% ============================================================
%  CHAPTER 6: CIRCULAR MOTION
%  The Physics Codex: A Complete University Foundation
%  Korede Samuel Omotosho (Falcon)
%  Aurikrex Learning Systems, 2026
%
%  File: chapters/part1/ch06_circular_motion.tex
% ============================================================

\chapter{Circular Motion}
\label{ch:circular_motion}

\chapterquote{The most beautiful thing we can experience
is the mysterious. It is the source of all true art
and science.}{Albert Einstein}

\begin{objectivebox}
By the end of this chapter, you should be able to:
\begin{enumerate}[noitemsep]
  \item Define angular displacement, angular velocity
  and angular acceleration
  \item Relate linear and angular quantities
  \item Explain why uniform circular motion requires
  a centripetal acceleration and force
  \item Calculate centripetal acceleration and
  centripetal force
  \item Apply circular motion equations to banked
  roads, conical pendulums, and vertical circles
  \item Explain the concept of centrifugal force as
  a pseudo-force
  \item Solve problems involving objects moving in
  vertical circles
\end{enumerate}
\end{objectivebox}

\vspace{0.4cm}

\minitoc

\vspace{0.5cm}

% ============================================================
\section{Why Circular Motion is Accelerated Motion}
% ============================================================

Here is a fact that surprises many students: a car
travelling around a roundabout at perfectly constant
speed is \textit{accelerating}. This seems contradictory
--- acceleration means changing velocity, and if the
speed is constant, how can velocity be changing?

The answer lies in the vector nature of velocity.
\textbf{Velocity is not just speed --- it is speed in
a specific direction.} As the car goes around the
roundabout, its direction changes continuously. A
changing direction means a changing velocity vector,
and a changing velocity means acceleration. This
acceleration is called \textbf{centripetal acceleration},
and it always points toward the centre of the circle.

\begin{misconceptionbox}
Many students confuse centripetal force with a real
outward force on the rotating object. There is no
outward force on a car going around a bend. The car
is pulled \textit{inward} by friction. What you feel
pressed against the car door is your own body's
inertia (Newton's First Law) --- your body wants to
continue in a straight line, but the car door forces
you to turn. This apparent outward force is called
\textbf{centrifugal force} and it is a \textit{pseudo-force}
--- it only appears to exist in the rotating reference
frame.
\end{misconceptionbox}

% ============================================================
\section{Angular Quantities}
% ============================================================

\begin{definitionbox}[Angular Displacement]
The \textbf{angular displacement} $\theta$ is the angle
swept out by the radius vector as the object moves along
the circular path. Measured in \textbf{radians} (rad).

One full circle $= 360° = 2\pi\ \text{rad}$\\
$1\ \text{rad} = 180°/\pi \approx 57.3°$
\end{definitionbox}

\begin{definitionbox}[Angular Velocity]
The \textbf{angular velocity} $\omega$ is the rate of
change of angular displacement:
\[
  \omega = \frac{\Delta\theta}{\Delta t}
\]
SI unit: rad\,s$^{-1}$.

For uniform circular motion (period $T$, frequency $f$):
\[
  \omega = \frac{2\pi}{T} = 2\pi f
\]
\end{definitionbox}

\begin{definitionbox}[Relationship Between Linear and
Angular Quantities]
For a body moving in a circle of radius $r$:
\[
  s = r\theta \quad \text{(arc length)}
\]
\[
  v = r\omega \quad \text{(linear speed)}
\]
\[
  a_{tangential} = r\alpha \quad
  \text{(tangential acceleration, if } \omega
  \text{ is changing)}
\]
where $\alpha = d\omega/dt$ is the angular acceleration.
\end{definitionbox}

\begin{diagbox}[Angular and Linear Quantities]
\begin{center}
\begin{tikzpicture}[scale=1.2]
  % Circle
  \draw[thick, codexblue!40] (0,0) circle (2cm);

  % Radius to position 1
  \draw[thick, codexblue] (0,0) -- (2,0)
    node[right]{$P_1$};
  \fill[codexblue] (2,0) circle (2.5pt);

  % Radius to position 2
  \draw[thick, codexred] (0,0) -- (1.41,1.41)
    node[above right]{$P_2$};
  \fill[codexred] (1.41,1.41) circle (2.5pt);

  % Angle
  \draw[codexgold, thick] (0.6,0) arc(0:45:0.6);
  \node[codexgold] at (0.85,0.25) {$\theta$};

  % Arc
  \draw[codexgreen, very thick]
    (2,0) arc(0:45:2);
  \node[codexgreen] at (2.3,0.8) {$s = r\theta$};

  % Radius label
  \node[below] at (1,-0.1) {$r$};

  % Velocity vector
  \draw[->, very thick, codexred]
    (1.41,1.41) -- ++(-1.0, 1.0)
    node[above left]{$v = r\omega$};

  % Centre
  \fill (0,0) circle (2pt) node[below left]{O};
\end{tikzpicture}
\end{center}
\end{diagbox}

\begin{examplebox}[6.1]
\stepcounter{workedexample}
A wheel of radius $0.40\ \text{m}$ rotates at
$300\ \text{rpm}$ (revolutions per minute).
Find: (a) the angular velocity in rad\,s$^{-1}$;
(b) the linear speed of a point on the rim;
(c) the period of rotation.
\end{examplebox}

\begin{solutionbox}
\textbf{(a)} Convert rpm to rad\,s$^{-1}$:
\[
  \omega = 300\ \frac{\text{rev}}{\text{min}}
  \times \frac{2\pi\ \text{rad}}{\text{rev}}
  \times \frac{1\ \text{min}}{60\ \text{s}}
  = \frac{300 \times 2\pi}{60}
  = 10\pi \approx 31.4\ \text{rad\,s}^{-1}
\]

\textbf{(b)} Linear speed of rim:
\[
  v = r\omega = 0.40 \times 10\pi
  = 4\pi \approx 12.6\ \text{m\,s}^{-1}
\]

\textbf{(c)} Period:
\[
  T = \frac{2\pi}{\omega} = \frac{2\pi}{10\pi}
  = 0.20\ \text{s}
\]
\end{solutionbox}

% ============================================================
\section{Centripetal Acceleration and Force}
% ============================================================

\begin{lawbox}[Centripetal Acceleration]
For an object moving at constant speed $v$ in a
circle of radius $r$, the centripetal acceleration
always points toward the centre of the circle:
\[
  a_c = \frac{v^2}{r} = \omega^2 r = \frac{v^2}{r}
\]
This is not an additional force --- it is the
\textit{resultant} of all real forces on the object,
directed toward the centre.
\end{lawbox}

\begin{lawbox}[Centripetal Force]
The \textbf{centripetal force} is the net force that
produces centripetal acceleration:
\[
  F_c = ma_c = \frac{mv^2}{r} = m\omega^2 r
\]
It always points toward the centre of the circle.
It is \textbf{not} a new type of force --- it is
provided by whatever real force is acting: tension,
gravity, friction, normal force, or a combination.
\end{lawbox}

\begin{center}
\begin{tabular}{lll}
\toprule
\textbf{Situation} & \textbf{What provides centripetal force}\\
\midrule
Ball on a string (horizontal) & Tension in the string\\
Planet orbiting the Sun & Gravitational force of the Sun\\
Car turning a corner & Friction between tyres and road\\
Electron orbiting nucleus & Electrostatic attraction\\
Clothes in washing machine & Normal force from drum wall\\
Pilot banking an aircraft & Component of lift force\\
\bottomrule
\end{tabular}
\end{center}

\begin{examplebox}[6.2]
\stepcounter{workedexample}
A $0.50\ \text{kg}$ ball is attached to a string of
length $1.2\ \text{m}$ and swings in a horizontal
circle. The string makes an angle of $30°$ with the
vertical. Find: (a) the tension in the string;
(b) the speed of the ball; (c) the period of revolution.
($g = 10\ \text{m\,s}^{-2}$)
\end{examplebox}

\begin{solutionbox}
This is a conical pendulum. The ball moves in a
horizontal circle of radius $r = L\sin\theta$.

\begin{center}
\begin{tikzpicture}[scale=1.0]
  % String
  \draw[thick] (0,3) -- (1.2*0.5,3-1.2*0.866);
  \fill (0,3) circle (2.5pt);

  % Ball
  \fill[codexblue!50] (0.6,3-1.039) circle (5pt)
    node[right, font=\small]{$m$};

  % Forces
  \draw[->, very thick, codexred]
    (0.6,3-1.039) -- (0,3)
    node[above left]{$T$};
  \draw[->, very thick, codexblue]
    (0.6,3-1.039) -- (0.6, 3-1.039-1.2)
    node[below]{$mg$};
  \draw[->, very thick, codexgreen]
    (0.6,3-1.039) -- (0.6-1.0, 3-1.039)
    node[left]{$F_c$};

  % Angle
  \draw[dashed, gray] (0,3) -- (0,1.5);
  \draw[codexgold] (0,2.7) arc(270:300:0.3);
  \node[codexgold, right, font=\small] at (0.15,2.55)
    {$30°$};

  % Radius: horizontal distance from the rotation axis to the bob.
  \draw[densely dotted, gray] (0,1.961) -- (0,1.7);
  \draw[densely dotted, gray] (0.6,1.961) -- (0.6,1.7);
  \draw[<->, dashed, gray] (0,1.7) -- (0.6,1.7);
  \node[below, gray, font=\small] at (0.3,1.7)
    {$r$};
\end{tikzpicture}
\end{center}

Resolving forces on the ball:

\textbf{Vertical equilibrium:}
\[
  T\cos 30° = mg
\]
\[
  T = \frac{mg}{\cos 30°}
  = \frac{0.50 \times 10}{0.866}
  = 5.77\ \text{N}
\]

\textbf{Horizontal (centripetal direction):}
\[
  T\sin 30° = \frac{mv^2}{r}
  = m\omega^2 r
\]
\[
  r = L\sin 30° = 1.2 \times 0.5 = 0.60\ \text{m}
\]
\[
  v^2 = \frac{T\sin 30° \cdot r}{m}
  = \frac{5.77 \times 0.5 \times 0.60}{0.50}
  = 3.46\ \text{m}^2\text{s}^{-2}
\]
\[
  v = 1.86\ \text{m\,s}^{-1}
\]

\textbf{Period:}
\[
  T_{period} = \frac{2\pi r}{v}
  = \frac{2\pi \times 0.60}{1.86}
  = 2.03\ \text{s}
\]
\end{solutionbox}

% ============================================================
\section{Motion on Banked Roads and Curves}
% ============================================================

When a vehicle goes around a curve, the required
centripetal force is provided by friction (on a flat
road) or by a component of the normal force (on a
banked road). Banking allows higher speeds without
relying on friction.

\begin{lawbox}[Ideal Banking Angle]
For a road banked at angle $\theta$, the ideal speed
$v$ (no friction needed) for a curve of radius $r$ is:
\[
  \tan\theta = \frac{v^2}{rg}
  \implies v = \sqrt{rg\tan\theta}
\]
\end{lawbox}

\begin{examplebox}[6.3]
\stepcounter{workedexample}
A racing circuit has a banked curve of radius
$200\ \text{m}$ banked at $18°$. Find the ideal
speed for this curve. ($g = 10\ \text{m\,s}^{-2}$,
$\tan 18° = 0.3249$)
\end{examplebox}

\begin{solutionbox}
\[
  v = \sqrt{rg\tan\theta}
  = \sqrt{200 \times 10 \times 0.3249}
  = \sqrt{649.8}
  = 25.5\ \text{m\,s}^{-1}
  \approx 91.8\ \text{km\,h}^{-1}
\]
At this speed, no friction is needed on the banked road.
\end{solutionbox}

\begin{realworldbox}[Lagos--Ibadan Expressway Curves]
Highway engineers in Nigeria must account for the
banking angle, speed limit, and coefficient of friction
when designing curves. The Lagos--Ibadan Expressway
has several curves where inadequate banking combined
with high speeds causes accidents, especially when
roads are wet and friction is reduced. Understanding
circular motion directly informs road safety design.
\end{realworldbox}

% ============================================================
\section{Motion in a Vertical Circle}
% ============================================================

When an object moves in a vertical circle, both gravity
and the centripetal force requirement change with
position, making the analysis richer than the
horizontal case.

\begin{processbox}[Object on a String in a Vertical Circle]
Consider a ball of mass $m$ on a string of length $r$
moving in a vertical circle. Let $v$ be the speed at
the bottom and $v'$ the speed at the top.

\textbf{At the BOTTOM of the circle:}
The tension $T_B$ must support the weight AND provide
centripetal force (both point in the same direction
--- toward the centre which is above):
\[
  T_B - mg = \frac{mv_B^2}{r}
  \implies T_B = mg + \frac{mv_B^2}{r}
\]

\textbf{At the TOP of the circle:}
Both tension and weight point downward (toward centre):
\[
  T_T + mg = \frac{mv_T^2}{r}
  \implies T_T = \frac{mv_T^2}{r} - mg
\]

\textbf{Minimum speed at the top} (string just stays
taut, $T_T = 0$):
\[
  0 = \frac{mv_{T,min}^2}{r} - mg
  \implies v_{T,min} = \sqrt{gr}
\]
\end{processbox}

\begin{diagbox}[Vertical Circle --- Forces at Key Points]
\begin{center}
\begin{tikzpicture}[scale=1.1]
  % Circle
  \draw[thick, gray!50] (0,0) circle (2cm);
  \fill (0,0) circle (2pt) node[right, font=\small]{O};

  % Bottom point
  \fill[codexblue] (0,-2) circle (4pt);
  \node[below, font=\small] at (-0.75,-2) {Bottom};
  \draw[->, very thick, codexred]
    (0,-2) -- (0,-0.7) node[right, font=\small]{$T_B$};
  \draw[->, very thick, codexblue]
    (0,-2) -- (0,-2-0.8) node[below, font=\small]{$mg$};

  % Top point
  \fill[codexgreen] (0,2) circle (4pt)
    node[above]{Top};
  % The two inward forces are offset sideways so both arrows remain visible.
  \draw[->, very thick, codexred]
    (-0.18,2) -- (-0.18,1.0) node[left, font=\small]{$T_T$};
  \draw[->, very thick, codexblue]
    (0.18,2) -- (0.18,1.5) node[right, font=\small]{$mg$};

  % Speed vectors
  \draw[->, very thick, codexgold]
    (0,-2) -- (1.2,-2)
    node[right, font=\small]{$v_B$ (max speed)};
  \draw[->, very thick, codexgold]
    (0,2) -- (-1.2,2)
    node[left, font=\small]{$v_T$ (min speed)};

  % Note
  \node[font=\tiny, align=center] at (0,-3.35)
    {At bottom: $T_B > mg$ (feels heavier)\\
    At top: $T_T + mg$ both centripetal};
\end{tikzpicture}
\end{center}
\end{diagbox}

\begin{examplebox}[6.4]
\stepcounter{workedexample}
A $0.40\ \text{kg}$ ball on a $0.80\ \text{m}$ string
moves in a vertical circle. The speed at the bottom
is $5.0\ \text{m\,s}^{-1}$.\\
(a) Find the tension at the bottom.\\
(b) Find the speed at the top using energy conservation.\\
(c) Find the tension at the top.\\
(d) Can the ball complete the full circle?
($g = 10\ \text{m\,s}^{-2}$)
\end{examplebox}

\begin{solutionbox}
\textbf{(a)} Tension at the bottom:
\[
  T_B = mg + \frac{mv_B^2}{r}
  = 0.40(10) + \frac{0.40(5.0)^2}{0.80}
  = 4.0 + 12.5 = 16.5\ \text{N}
\]

\textbf{(b)} Speed at top (energy conservation,
height difference $= 2r = 1.6\ \text{m}$):
\[
  \frac{1}{2}mv_T^2 = \frac{1}{2}mv_B^2 - mg(2r)
\]
\[
  v_T^2 = v_B^2 - 4gr
  = 25 - 4(10)(0.80) = 25 - 32 = -7
\]

$v_T^2$ is negative --- this is impossible.
The ball \textbf{cannot} reach the top.

\textbf{Alternative check using minimum speed:}
\[
  v_{T,min} = \sqrt{gr} = \sqrt{10 \times 0.80}
  = \sqrt{8} = 2.83\ \text{m\,s}^{-1}
\]
Required speed at top would need:
\[
  v_B^2 \geq v_{T,min}^2 + 4gr
  = 8 + 32 = 40
  \implies v_B \geq \sqrt{40} = 6.32\ \text{m\,s}^{-1}
\]
Since $v_B = 5.0\ \text{m\,s}^{-1} < 6.32\ \text{m\,s}^{-1}$,
the ball cannot complete the circle. The string
goes slack before the top is reached.
\end{solutionbox}

% ============================================================
\section{Satellites and Orbital Motion}
% ============================================================

A satellite in a circular orbit is in a state of
continuous free fall --- it is falling toward the
Earth, but its horizontal speed is so high that the
Earth curves away beneath it at the same rate it
falls. Gravity provides the centripetal force.

\begin{lawbox}[Orbital Speed and Period]
For a satellite orbiting at radius $r$ from Earth's
centre (mass $M_E$):
\[
  \frac{GMm}{r^2} = \frac{mv^2}{r}
  \implies v = \sqrt{\frac{GM}{r}}
\]
\[
  T = \frac{2\pi r}{v} = 2\pi r\sqrt{\frac{r}{GM}}
  = 2\pi\sqrt{\frac{r^3}{GM}}
\]
This is Kepler's Third Law: $T^2 \propto r^3$.
\end{lawbox}

\begin{examplebox}[6.5]
\stepcounter{workedexample}
A satellite orbits Earth at a height of
$400\ \text{km}$ above the surface. Find the orbital
speed and period. ($R_E = 6.4 \times 10^6\ \text{m}$,
$g_s = 9.81\ \text{m\,s}^{-2}$,
$GM = gR_E^2 = 9.81 \times (6.4\times10^6)^2$)
\end{examplebox}

\begin{solutionbox}
Orbital radius:
\[
  r = R_E + h = 6.4\times10^6 + 4\times10^5
  = 6.8\times10^6\ \text{m}
\]

$GM = g_s R_E^2 = 9.81 \times (6.4\times10^6)^2
= 4.016\times10^{14}\ \text{m}^3\text{s}^{-2}$

Orbital speed:
\[
  v = \sqrt{\frac{GM}{r}}
  = \sqrt{\frac{4.016\times10^{14}}{6.8\times10^6}}
  = \sqrt{5.906\times10^7}
  = 7685\ \text{m\,s}^{-1}
  \approx 7.7\ \text{km\,s}^{-1}
\]

Period:
\[
  T = \frac{2\pi r}{v}
  = \frac{2\pi \times 6.8\times10^6}{7685}
  = 5556\ \text{s}
  \approx 92.6\ \text{min}
\]

This matches the International Space Station's
actual orbital period of approximately 92 minutes.
\end{solutionbox}

% ============================================================
\section{Chapter Summary}
% ============================================================

\begin{summarybox}
\textbf{Angular velocity:}
$\omega = \Delta\theta/\Delta t = 2\pi f = 2\pi/T$
(rad\,s$^{-1}$)

\textbf{Linear--angular relations:}
$v = r\omega$, $s = r\theta$

\textbf{Centripetal acceleration:}
$a_c = v^2/r = \omega^2 r$ (toward centre)

\textbf{Centripetal force:}
$F_c = mv^2/r = m\omega^2 r$ (net inward force)

\textbf{Centrifugal force:} pseudo-force in rotating
frame only --- not a real force.

\textbf{Conical pendulum:}
$\tan\theta = v^2/(rg)$,
$T\cos\theta = mg$, $T\sin\theta = mv^2/r$

\textbf{Banking angle:}
$\tan\theta = v^2/(rg)$ (ideal, no friction)

\textbf{Vertical circle:}
$T_{bottom} = mg + mv^2/r$;
$T_{top} = mv^2/r - mg$;
minimum speed at top $= \sqrt{gr}$

\textbf{Orbital speed:}
$v = \sqrt{GM/r}$;
$T^2 \propto r^3$ (Kepler's Third Law)
\end{summarybox}

% ============================================================
\newpage
\begin{exercisebox}[6]

\subsection*{Section A: Multiple Choice}

\textbf{A1.} A body moves at constant speed in a
circle. Which of the following is true?

\mcoption{A}{Its velocity and acceleration are both
constant}
\mcoption{B}{Its speed is constant but velocity
changes direction}
\mcoption{C}{No force acts on it}
\mcoption{D}{Its kinetic energy continuously increases}
\ansref{Ch6-A1}

\vspace{0.4cm}

\textbf{A2.} A stone of mass $0.2\ \text{kg}$ is
whirled in a horizontal circle of radius $0.5\ \text{m}$
at $4\ \text{m\,s}^{-1}$. The centripetal force is:

\mcoption{A}{$1.6\ \text{N}$}
\mcoption{B}{$3.2\ \text{N}$}
\mcoption{C}{$6.4\ \text{N}$}
\mcoption{D}{$12.8\ \text{N}$}
\ansref{Ch6-A2}

\vspace{0.4cm}

\textbf{A3.} A satellite orbits at radius $r$ from
Earth's centre with period $T$. If it moves to orbit
at radius $4r$, its new period is:

\mcoption{A}{$2T$}
\mcoption{B}{$4T$}
\mcoption{C}{$8T$}
\mcoption{D}{$16T$}
\ansref{Ch6-A3}

\vspace{0.4cm}

\textbf{A4.} A car moves around a flat circular bend
of radius $50\ \text{m}$. The maximum speed before
sliding occurs is (coefficient of friction $= 0.4$,
$g = 10\ \text{m\,s}^{-2}$):

\mcoption{A}{$10\ \text{m\,s}^{-1}$}
\mcoption{B}{$14.1\ \text{m\,s}^{-1}$}
\mcoption{C}{$20\ \text{m\,s}^{-1}$}
\mcoption{D}{$200\ \text{m\,s}^{-1}$}
\ansref{Ch6-A4}

\vspace{0.4cm}

\textbf{A5.} The minimum speed at the top of a
vertical circle of radius $r$ for a ball on a
string to maintain tension is:

\mcoption{A}{$\sqrt{rg/2}$}
\mcoption{B}{$\sqrt{rg}$}
\mcoption{C}{$\sqrt{2rg}$}
\mcoption{D}{$\sqrt{5rg}$}
\ansref{Ch6-A5}

\vspace{0.4cm}

\textbf{A6.} A wheel rotates at $10\ \text{rad\,s}^{-1}$.
The angular displacement in $5\ \text{s}$ is:

\mcoption{A}{$2\ \text{rad}$}
\mcoption{B}{$10\ \text{rad}$}
\mcoption{C}{$50\ \text{rad}$}
\mcoption{D}{$100\ \text{rad}$}
\ansref{Ch6-A6}

\vspace{0.4cm}

\textbf{A7.} For circular motion on a banked road
at the correct (ideal) speed, the frictional force
on the car is:

\mcoption{A}{Directed toward the centre}
\mcoption{B}{Directed away from the centre}
\mcoption{C}{Zero}
\mcoption{D}{Equal to the normal reaction}
\ansref{Ch6-A7}

\vspace{0.4cm}

\textbf{A8.} The centripetal acceleration of the Moon
in its orbit is approximately $0.0027\ \text{m\,s}^{-2}$.
The distance from the Earth to the Moon is
$3.84 \times 10^8\ \text{m}$. The orbital speed of the
Moon is approximately:

\mcoption{A}{$720\ \text{m\,s}^{-1}$}
\mcoption{B}{$1020\ \text{m\,s}^{-1}$}
\mcoption{C}{$1560\ \text{m\,s}^{-1}$}
\mcoption{D}{$2440\ \text{m\,s}^{-1}$}
\ansref{Ch6-A8}

\vspace{0.4cm}

\textbf{A9.} A car of mass $1000\ \text{kg}$ travels
over a hill that has a circular cross-section of
radius $40\ \text{m}$. The maximum speed at the top
of the hill without leaving the road is
($g = 10\ \text{m\,s}^{-2}$):

\mcoption{A}{$10\ \text{m\,s}^{-1}$}
\mcoption{B}{$20\ \text{m\,s}^{-1}$}
\mcoption{C}{$40\ \text{m\,s}^{-1}$}
\mcoption{D}{$400\ \text{m\,s}^{-1}$}
\ansref{Ch6-A9}

\vspace{0.4cm}

\textbf{A10.} A ball on a string of length $L$ is
swung in a vertical circle. At the top of the circle,
both gravity and tension act downward. If the tension
at the top is equal to the weight $mg$, the speed
at the top is:

\mcoption{A}{$\sqrt{gL}$}
\mcoption{B}{$\sqrt{2gL}$}
\mcoption{C}{$\sqrt{gL/2}$}
\mcoption{D}{$2\sqrt{gL}$}
\ansref{Ch6-A10}

\vspace{0.6cm}

\subsection*{Section B: Theory and Calculations}

\textbf{B1.} A $0.30\ \text{kg}$ ball is attached
to a string of length $0.60\ \text{m}$ and swings
as a conical pendulum, making an angle of $45°$ with
the vertical. ($g = 10\ \text{m\,s}^{-2}$,
$\tan 45° = 1$, $\cos 45° = 0.707$)\\[4pt]
(a) Calculate the tension in the string.\\
(b) Calculate the radius of the circular path.\\
(c) Calculate the speed of the ball.\\
(d) Calculate the period of revolution.\\
(e) If the string can withstand a maximum tension
of $10\ \text{N}$, what is the maximum angle
$\theta$ before the string breaks?
\hfill\ansref{Ch6-B1}

\vspace{0.4cm}

\textbf{B2.} A car of mass $1200\ \text{kg}$ travels
around a circular track of radius $150\ \text{m}$.
The track is banked at $12°$ and the coefficient of
static friction is $0.35$. ($g = 10\ \text{m\,s}^{-2}$,
$\sin 12° = 0.208$, $\cos 12° = 0.978$,
$\tan 12° = 0.213$)\\[4pt]
(a) Find the ideal speed for the banked track
(no friction needed).\\
(b) Find the maximum safe speed on this track.
(Hint: at maximum speed, friction acts down the slope
to prevent the car sliding outward.)\\
(c) Find the minimum safe speed
(friction acts up the slope).
\hfill\ansref{Ch6-B2}

\vspace{0.4cm}

\textbf{B3.} A $0.25\ \text{kg}$ ball on a
$1.5\ \text{m}$ string moves in a vertical circle.
($g = 10\ \text{m\,s}^{-2}$)\\[4pt]
(a) Find the minimum speed at the top for the
string to remain taut.\\
(b) Using energy conservation, find the minimum
speed at the bottom for the ball to complete
the full circle.\\
(c) At this minimum bottom speed, calculate the
tension at the bottom.\\
(d) Calculate the tension at a point where the
string is horizontal (side), using energy
conservation to find the speed there.
\hfill\ansref{Ch6-B3}

\vspace{0.4cm}

\textbf{B4.} A geostationary satellite orbits Earth
with a period equal to Earth's rotation period
(24 hours). ($G = 6.67\times10^{-11}\ \text{N\,m}^2\text{kg}^{-2}$,
$M_E = 5.97\times10^{24}\ \text{kg}$)\\[4pt]
(a) Show that the orbital radius of a geostationary
satellite is given by
$r = \left(\frac{GM_E T^2}{4\pi^2}\right)^{1/3}$.\\
(b) Calculate the orbital radius.\\
(c) Calculate the orbital speed.\\
(d) Explain why a geostationary satellite must orbit
above the equator and why this is useful for
telecommunications.
\hfill\ansref{Ch6-B4}

\vspace{0.4cm}

\textbf{B5.} In a theme park ride, passengers sit in
cars that travel in a vertical circle of radius
$12\ \text{m}$. The cars complete one revolution
every $6.0\ \text{s}$.
($g = 10\ \text{m\,s}^{-2}$, $\pi^2 \approx 10$)\\[4pt]
(a) Calculate the angular velocity of the cars.\\
(b) Calculate the linear speed.\\
(c) Calculate the centripetal acceleration.\\
(d) A $70\ \text{kg}$ passenger sits in the car.
Find the normal force (apparent weight) at the
bottom and at the top of the circle.\\
(e) At the top, does the passenger feel heavier
or lighter than normal? Explain using Newton's
Second Law.
\hfill\ansref{Ch6-B5}

\end{exercisebox}
